\input{../tex-inputs/latex-leseansicht-vorspann}

\section[Arthur Schnitzler an Theodor von Sosnosky, 30. 12. 1899]{L04236 Arthur Schnitzler an Theodor von Sosnosky, 30. 12. 1899}
\nopagebreak\mylabel{L04236v}
\rehead{ }\normalsize\beginnumbering\briefempfaengerindex{Sosnosky, Theodor von@\textsc{Sosnosky, Theodor von}!zzzSchnitzler, Arthur@\emph{von Arthur Schnitzler}!1899-12-302@{30. 12. 1899}|(be}
\toendnotes[C]{\smallbreak\pagebreak[2]}
\correspDesc{Versand  durch Arthur Schnitzler am 30. 12. 1899 in Wien
\newline{}Erhalt  durch Theodor von Sosnosky im Zeitraum [31. 12. 1899 – 4. 1. 1900?] \textbf{Ort fehlend} }\toendnotes[C]{\smallbreak}
\buchAlsQuelle{Theodor v. Sosnosky: \emph{Artur Schnitzler als Lyriker.} In: \emph{Neues Wiener Journal}, Jg. 40, Nr. 13.823, 18. 5. 1932, S. 6–7.}\toendnotes[C]{\smallbreak}
\pstart{}{\pb}Sehr geehrter Herr v. Sosnosky!\pend\vspace{0.5em}
\pstart
           In einigen Tagen werde ich Ihnen ein paar Gedichte\pwindex{Schnitzler, Arthur 15.\,5.\,1862 Wien – 21.\,10.\,1931 ebd.@\textsc{Schnitzler, Arthur} (15.\,5.\,1862 Wien – 21.\,10.\,1931 ebd.), \emph{Schriftsteller, Mediziner}!Ohnmacht@\strich\emph{Ohnmacht}|pwv}\pwindex{Schnitzler, Arthur 15.\,5.\,1862 Wien – 21.\,10.\,1931 ebd.@\textsc{Schnitzler, Arthur} (15.\,5.\,1862 Wien – 21.\,10.\,1931 ebd.), \emph{Schriftsteller, Mediziner}!Anfang vom Ende@\strich\emph{Anfang vom Ende}|pwv}\pwindex{Schnitzler, Arthur 15.\,5.\,1862 Wien – 21.\,10.\,1931 ebd.@\textsc{Schnitzler, Arthur} (15.\,5.\,1862 Wien – 21.\,10.\,1931 ebd.), \emph{Schriftsteller, Mediziner}!Wie wir so still@\strich\emph{Wie wir so still}|pwv} in Abſchrift{ }ſenden (ſo daß Sie{ }ſie nicht kopieren zu laſſen brauchen). Ich habe kaum ein Dutzend
               geſchrieben, die angehen. Was Sie an Lyrik an mir finden, hat{ }ſich in meine
               Erzählungen und Stücke hineingeſchlichen – was nicht immer von Vorteil war.\pend
           
\pstart
           Mit verbindlichen Grüßen{\\[\baselineskip]}hochachtungsvoll{\\[\baselineskip]}\spacefill\mbox{Dr. Artur Schnitzler.}\pend
           \leftskip=0em{}
\pstart
           30.\,12.\,99.\pend
           \selectlanguage{ngerman}\endnumbering\briefempfaengerindex{Sosnosky, Theodor von@\textsc{Sosnosky, Theodor von}!zzzSchnitzler, Arthur@\emph{von Arthur Schnitzler}!1899-12-302@{30. 12. 1899}|)be}\mylabel{L04236h}
\begin{anhang}
\end{anhang}\newcommand{\dateiname}{L04236}\newcommand{\titel}{Arthur Schnitzler an Theodor von Sosnosky, 30. 12. 1899}\newcommand{\editorInnen}{Herausgegeben von Jahnke, SelmaMüller, Martin Anton}\input{../tex-inputs/latex-leseansicht-abspann}
